\section*{Введение}
\addcontentsline{toc}{section}{Введение}

\textbf{Задание лабораторной работы~1.}
\begin{enumerate}
    \item Сформировать случайным образом связный ациклический граф в соответствии с заданным распределением (параметры распределения задаются как константы и подбираются экспериментально) с введенным пользователем количеством вершин. Распределение брать из справочника Вадзинского (есть в списке литературы): либо реализовать блок-схему, либо генерирующую формулу (в справочнике дан один из этих вариантов). Генерация происходит по степеням вершин. В зависимости от дальнейших заданий будет использовать ориентированный или неориентированный граф. Можно реализовывать в консоли.

    \item Посчитать экцентриситеты вершин, выделить центр графа и диаметральные вершины.

    \item Реализовать метод Шимбелла для сгенерированной с помощью заданного распределения весовой матрицы  (пользователь вводит количество ребер), в ответе представить минимальный и/или максмиальный путь  (в виде матрицы). Весовая матрица генерируются также в соотвествии с заданным распределением. При генерации необходимо учесть генерацию только положительных значений, только отрицательных значений и смешанную (на выбор пользователя).

    \item Определить возможность построения маршрута от одной заданной точки до другой (вершины вводит пользователь) и указывать количество таковых маршрутов.

    \item Реализовать меню для управления всеми пунктами.
\end{enumerate}

\textbf{Задание лабораторной работы~2.}
\begin{enumerate}
    \item Для заданных графов (случайно сгенерированных в предыдущей работе)
    выполнить обход рёбер графа поиском в ширину.
    \item Сгенерировать весовую матрицу (положительную или с отрицательными
    весами~--- выбирает пользователь) и найти кратчайший путь для двух
    выбранных пользователем вершин. В результате выводится не только
    расстояние, но и сам путь в виде последовательности вершин.
    Реализовать классический алгоритм Дейкстры с выводом вектора расстояний.
    \item Сравнить скорости работы реализованных алгоритмов по количеству
    итераций.
\end{enumerate}


Реализация выполнена на C++20. Проект разбит на статические библиотеки,
собираемые через CMake.